\clearpage\thispagestyle{empty}
\pdfbookmark[0]{Quellen zur Abhandlung über die Bestimmung der Polhöhe}{poli1808-sources}
\begin{center}\textbf{Quellen zur Abhandlung über die Bestimmung der Polhöhe}\end{center}
\noindent Diese KI-Übersetzung folgt dem lateinischen Text in den \textit{Werken}, Band~VI (Göttingen, 1874), S.~37--49. Sie enthält das Titelblatt, alle sieben Artikel, Gauß’ Anmerkung zu den Azimuten und beide Vorlesungsankündigungen. Auf dem Titelblatt steht \mbox{\textsc{elevationem}}, in der wiederholten Überschrift \mbox{\textsc{elpvationem}}. Beide Überschriften sind dem Sinn nach gleich übersetzt. \href{https://gdz.sub.uni-goettingen.de/id/PPN236007467}{Göttinger Digitalisat}.
\par\smallskip\noindent Zur astronomischen Terminologie siehe Siegfried Krauter, \textit{Mathematische Erd- und Himmelskunde}, Fassung vom Januar 2012, S.~24--26.
\par\smallskip\noindent Übersetzung und Satz wurden mit KI erstellt.
\clearpage
